\chapter{ЗАДАЧА ПРУЖНОСТІ У М'ЯКИХ ТКАНИНАХ ТІЛА ЛЮДИНИ}

%Hidden cite for correct bibliography ordering
\nocite{bahvalov-et-al,benerdge-et-al} 

\section{Огляд проблеми}

Одним із нововведень у малоінвазивній робототехніці був тактильний відгук від інструментів. Це була суттєва різниця між
відкритою операцією і малоінвазивною хірургією. Під час відкритої операції хірург міг\textbackslashмогла відчути певні
характеристики руками і оцінити властивості м'яких тканин. Тільки остання версія робота хірурга da Vinci 5 від
Intuitive Machines отримав таку функціональність \cite{elast-intiitive-machines-force-feedback}. Першим сертифікованим 
роботом із такою функціональністю був Senhance Robotic System. Тим не менше, наразі обмежене коло роботів хірургів надає
таку функціональність.

Для того щоб можна було розробляти такі системи управління роботами із тактильним відгуком потрібно створити симульоване 
середовище, яке буде надавати достовірні дані для калібрування інструментів. Його також можна використати для навчання
хірургів і виробляння звички ефективного управління такою системою. Однією із переваг такого керування роботом була б 
здатність відчувати неоднорідні включення у м'яких тканинах, які по аналогії із відкритою операцією можна використати 
для виявлення пухлин чи інших злоякісних утворень.

У цьому розділі розглянуто дві задачі. Перша плоска, а друга - просторова стаціонарна задача контакту кінця хірургічного 
інструменту із неоднорідним включенням у м'яких тканинах. Подається візуалізація відповідних сил.

\section{Плоска стаціонарна задача}

\subsection{Формулювання задачі}

Розглянемо область у якій є дві частини (рис. \ref{fig:elasticity_2d_domain}). $\Omega_1$ відповідає м'яким тканинам, а
$\Omega_3$ - це злоякісне утворення. Область $\Omega_2$ - це хірургічний інструмент, який контактує із м'якими тканинами
по межі $\Gamma_2$, тобто $\Gamma_2=\partial\Omega_2 \cap \partial\Omega_1$, 
$\Gamma_1 = \partial \Omega_1 \backslash \Gamma_2$.

\begin{figure}[ht!]
    \centering
    \begin{tikzpicture}
        \draw (5.7,3) rectangle (6.3,7);
        \draw (1,1) rectangle (11,3);
        \draw (6,2.5) circle (0.5);
        \draw[thick, ->] (5.3,6) -- (5.3,4);

        \draw[thick,->] (0,0) -- (12,0) node[anchor=north west] {x};
        \draw[thick,->] (0,0) -- (0,9) node[anchor=south east] {y};        
    
        \foreach \x in {0,...,4}
            \pgfmathtruncatemacro{\label}{\x * 10}
            \draw ( \x * 2.5 cm,1pt) -- ( \x * 2.5 cm,-1pt) node[anchor=north] {\label};
        \foreach \y in {0,...,3}
            \pgfmathtruncatemacro{\label}{\y * 10}
            \draw (1pt,\y * 2.5 cm) -- (-1pt,\y * 2.5 cm) node[anchor=east] {\label};
        
        \node at (11.5,2) {a};
        \node at (6,0.5) {b};
        \node at (6,7.5) {c};
        \node at (6,3.5) {$\Gamma_2$};

        \node at (8,2) {$\Omega_1$};
        \node at (6,5) {$\Omega_2$};
        \node at (6,2.5) {$\Omega_3$};
        \node at (4,0.5) {$\Gamma_1$};

    \end{tikzpicture}
    
    \caption{Переріз тканин і інструмента}
    \label{fig:elasticity_2d_domain}
\end{figure}

\noindent Параметр a рівний 8 см, b --- 40 см, а товщина хірургічного інструменту c --- 6 мм. Радіус кола $\Omega_3$, 
$r_c$ дорівнює 5 мм.

Скориставшись у рівняннях рівноваги законом Гука:

\begin{equation}
    \label{eqn:elasticity_hooke_law}
    \sigma_{ij}(x) = \lambda(x)\delta_{ij}\varepsilon_{kk}(x) + 2\mu(x)\varepsilon_{ij}(x)
\end{equation}

\noindent та співвідношеннями Коші:

\begin{equation}
    \label{eqn:elasticity_gauchy_tensor}
    \varepsilon_{ij}(x) = \frac{1}{2}(u_{i,j}(x) + u_{j,i}(x))
\end{equation}

\noindent одержимо систему диференційних рівнянь у переміщеннях:

\begin{equation}
    \label{eqn:elasticity_2d_eqn}
    \mu(x)u_{i,jj}(x) + (\lambda(x) + \mu(x))u_{j,ji} = 0, x \in \Omega_1
\end{equation}

\noindent де $\sigma_{ij}$ і $\varepsilon_{ij}$ --- компоненти відповідно тензора напружень та тензора деформацій; $u_i$
--- компоненти вектора переміщень; $\delta_{ij}$ --- символ Кронекера; $\lambda$ і $\mu$ --- коефіцієнти Ляме.

\noindent На всій межі $\Gamma_1$ немає переміщень
\begin{equation}
    \label{eqn:elasticity_2d_cond_1}
    u_i(x) = 0, x \in \Gamma_1
\end{equation}

\noindent На всій межі $\Gamma_2$ діє хірургічний інструмент вздовж нормалі із компонентами $n_j$ і зусиллям
\begin{equation}
    \label{eqn:elasticity_2d_cond_2}
    t_i(x) = t_{pi}, x \in \Gamma_2
\end{equation}

\noindent де $t_i = \sigma_{ij}n_j$.

\textbf{Коефіцієнти Ляме.} Для обчислення коефіцієнтів $\lambda$ і $\mu$ у здорових і уражених тканинах використовується
модуль Юнга і коефіцієнт Пуассона. Зокрема

\begin{equation}
    \label{eqn:elasticity_2d_lame_lambda_convertion}
    \lambda = \frac{E\nu}{(1+\nu)(1-2\nu)}
\end{equation}

\begin{equation}
    \label{eqn:elasticity_2d_lame_mu_convertion}
    \mu = \frac{E\nu}{2(1+\nu)}
\end{equation}

\noindent де $E$ --- модуль Юнга, а $\nu$ --- коефіцієнт Пуассона.  

На основі статті \cite{elast-young-modulus-cancer} візьмемо модулем Юнга для здорової тканини рівним 1060 Па, а у 
злоякісному утворенні він може досягати 8640 Па. У статті \cite{elast-poisson-rato-cancer} можна знайти відомості про 
коефіцієнт Пуассона. Для здорової тканини він має значення 0,31, а для ураженої зменшується до 0,26. 

Обчислимо коефіцієнти Ляме за формулами \ref{eqn:elasticity_2d_lame_lambda_convertion} і 
\ref{eqn:elasticity_2d_lame_mu_convertion}. Враховуючи неоднорідності отримаємо наступне представлення:

\begin{equation}
    \label{eqn:elasticity_2d_lyame_lambda}
    \lambda(x) = \lambda_c(x)\chi_3 + \lambda_t, x \in \Omega_1
\end{equation}

\begin{equation*}
    \lambda_c(x) = (3714,3 - \lambda_t)cos\left[\frac{\pi}{2}\frac{||x-x_c||^2}{r_c^2}\right], \lambda_t = 660,1
\end{equation*}

\begin{equation}
    \label{eqn:elasticity_2d_lyame_mu}
    \mu(x) = \mu_c(x)\chi_3 + \mu_t, x \in \Omega_1
\end{equation}

\begin{equation*}
    \mu_c(x) = (3428,6 - \mu_t)cos\left[\frac{\pi}{2}\frac{||x-x_c||^2}{r_c^2}\right], \mu_t = 404,6
\end{equation*}

\noindent де $\chi_3$ --- характеристична функція області $\Omega_3$, a $x_c$ --- центр області.

Рівняння \ref{eqn:elasticity_2d_eqn} перепишемо у наступному вигляді:

\begin{equation}
    \label{eqn:elasticity_2d_eqn_with_f}
    \mu_t u_{i,jj}(x) + (\lambda_t + \mu_t)u_{j,ji}(x) = -f_{ij,j}(\varepsilon(x)), x \in \Omega_1
\end{equation}

\noindent, де

\begin{equation*}
    f_{ij}(\varepsilon) = [ 2\mu_c(x)\varepsilon_{ij}(x) + \lambda_c(x) \delta_{ij} \varepsilon_{kk}(x) ] \chi_3
\end{equation*}

Отже, для визначення переміщень маємо змішану крайову задачу \ref{eqn:elasticity_2d_eqn_with_f},
\ref{eqn:elasticity_2d_cond_1} і \ref{eqn:elasticity_2d_cond_2} із коефіцієнтами Ляме у вигляді
\ref{eqn:elasticity_2d_lyame_lambda} і \ref{eqn:elasticity_2d_lyame_mu}.

\subsection{Дискретизація рівнянь}

\textbf{Інтегральні зображення розв'язків.} На основі прямого методу граничних елементів, як показано у
\cite{brebbia-bem}, опишемо розв'язок рівняння пружності (\ref{eqn:elasticity_2d_eqn_with_f}) для $\xi \in \Omega_1$, де
$\Gamma = \partial\Omega_1$:

\begin{equation}
    \label{eqn:elasticity_2d_bem_direct_elasticity}
    \begin{split}
    u_i(\xi) = \int_{\Gamma}\hat{u}_{ij}(x,\xi)p_j(x)d\Gamma(x) - \int_{\Gamma}\hat{p}_{ij}(x,\xi)u_j(x)d\Gamma(x) + \\  
        + \int_{\Omega_3} \hat{u}_{ij}(x,\xi)f_{jk,k}(\varepsilon(x))d\Omega(x)
    \end{split}
\end{equation}

\noindent$\hat{u}_{ij}$ і $\hat{p}_{ij}$ фундаментальні розв'язки плоского рівняння пружності для переміщень і напружень 
відповідно. Застосуємо до \ref{eqn:elasticity_2d_bem_direct_elasticity} теорему Гріна для інтегрування частинами. 
Одержимо:

\begin{equation}
    \label{eqn:elasticity_2d_bem_direct_elasticity_simplified}
    \begin{split}
    u_i(\xi) = \int_{\Gamma}\hat{u}_{ij}(x,\xi)p_j(x)d\Gamma(x) - \int_{\Gamma}\hat{p}_{ij}(x,\xi)u_j(x)d\Gamma(x) + \\ 
        + \int_{\Gamma_3} \hat{u}_{ij}(x,\xi)f_{jk}(\varepsilon(x))n_kd\Gamma(x)
        - \int_{\Omega_3} \hat{u}_{ij,k}(x,\xi)f_{jk}(\varepsilon(x))d\Omega(x)
    \end{split}
\end{equation}

Застосуємо закон Гука \ref{eqn:elasticity_hooke_law} і співвідношеннями Коші \ref{eqn:elasticity_gauchy_tensor} щоб 
отримати інтегральні представлення для тензора деформацій:

\begin{equation}
    \label{eqn:elasticity_2d_bem_direct_gauchy_tensor}
    \begin{split}
    \varepsilon_{ij}(\xi) = \int_{\Gamma}\tilde{u}_{ijk}(x,\xi)p_k(x)d\Gamma(x) 
        - \int_{\Gamma}\tilde{p}_{ijk}(x,\xi)u_k(x)d\Gamma(x) + \\ 
        + \int_{\Gamma_3} \tilde{u}_{ijk}(x,\xi)f_{kl}(\varepsilon(x))n_l d\Gamma(x)
        - \int_{\Omega_3} \tilde{u}_{ijk,l}(x,\xi)f_{kl}(\varepsilon(x))d\Omega(x),
    \end{split}
\end{equation}

\noindent де $\tilde{u}_{ijk} = \frac{1}{2} \{ \hat{u}_{ik,j} + \hat{u}_{jk,i} \} $ і 
$\tilde{p}_{ijk} = \frac{1}{2} \{ \hat{p}_{ik,j} + \hat{p}_{jk,i} \}$.

\textbf{Граничні і скінчені елементи.} Границі областей $\Omega_1$ і $\Omega_3$ розіб'ємо на граничні елементи. 
$\partial\Omega_1 = \cup_{v=1}^{V}\Gamma^v$, $\Gamma^n \cap \Gamma^m = \varnothing$, коли $n \neq m$. 
Додатково область $\Omega_3$ розіб'ємо на трикутні елементи. Частина із яких буде мати спільну сторону із 
граничними елементами, таким чином, що для будь якого граничного елемента знайдеться трикутник всередині $\Omega_3$, 
який містить цей елемент, як свою сторону. Тобто $\Omega_3 = \cup_{w=1}^{W} \Omega_3^w$, 
$\Omega_3^n \cap \Omega_3^m = \varnothing$, якщо $n \neq m$. 

\textbf{Система рівнянь.} На границі $\Gamma$ побудуємо систему апроксимаційних функцій і представимо переміщення і 
поверхневі зусилля у такому вигляді $\left.u(\eta_1)\right|_{\Gamma^m} = \phi^T(\eta_1)u^m$ і 
$\left.p(\eta_1)\right|_{\Gamma^n} = \phi^T(\eta_1)p^n$. Також на основі трикутних елементів і дотичних до них граничних 
елементів області $\Omega_3$ побудуємо такі функції для елементів тензора деформацій 
$\left.\varepsilon(\eta_1, \eta_2)\right|_{\Omega_3^l} = \psi^T(\eta_1,\eta_2)\varepsilon^l$. $\phi(\eta_1)$ і 
$\psi(\eta_1, \eta_2)$ --- інтерполяційні функції в локальних координатах для границі $\Gamma$ і області $\Omega_3$. 
Рівняння \ref{eqn:elasticity_2d_bem_direct_elasticity_simplified} і \ref{eqn:elasticity_2d_bem_direct_gauchy_tensor} 
представимо на основі інтерполяційних функцій наступним чином:

\begin{equation}
    \begin{split}
        \left\{ \begin{array}{lr} u_i \\ \varepsilon_{ij}\end{array} \right\}(x) =  
        \left\{ \begin{array}{lr} \hat{u}_{ij}^{\Gamma^v} \\ \tilde{u}_{ijk}^{\Gamma^v} \end{array} \right\}(x) p_j -
        \left\{ \begin{array}{lr} \hat{p}_{ij}^{\Gamma^v} \\ \tilde{p}_{ijk}^{\Gamma^v} \end{array} \right\}(x) u_j + \\
        + \left\{ \begin{array}{lr} \hat{u}_{ijk}^{\Gamma_3^w} \\ \tilde{u}_{ijkl}^{\Gamma_3^w} \end{array} \right\}(x) 
        \varepsilon_{jk}  
        - \left\{ \begin{array}{lr} \hat{u}_{ijk}^{\Omega_3^w} \\ \tilde{u}_{ijkl}^{\Omega_3^w} \end{array} \right\}(x) 
        \varepsilon_{jk},  
    \end{split}
\end{equation}

\noindent де

\begin{equation*}
    \left\{ \begin{array}{lr} \hat{u}_{ij}^{\Gamma^v} \\ \tilde{u}_{ijk}^{\Gamma^v} \end{array} \right\}(x) =
    \int_{-1}^{1} \left\{ \begin{array}{lr} \hat{u}_{ij} \\ \tilde{u}_{ijk} \end{array} \right\} 
    \phi^T J_{\Gamma_v} d \eta_1,
\end{equation*}

\begin{equation*}
    \left\{ \begin{array}{lr} \hat{p}_{ij}^{\Gamma^v} \\ \tilde{p}_{ijk}^{\Gamma^v} \end{array} \right\}(x) =
    \int_{-1}^{1} \left\{ \begin{array}{lr} \hat{p}_{ij} \\ \tilde{p}_{ijk} \end{array} \right\} 
    \phi^T J_{\Gamma_v} d \eta_1,
\end{equation*}

%TODO: Add more
\dots

Застосуємо метод колокацій по точках всередині границі і трикутних елементів в області неоднорідності $\Omega_3$. 
$\xi_0 \in \Gamma_v \cup \Omega_3$.

\begin{equation}
    \begin{split}
        \left\{ \begin{array}{lr} u_i \\ \varepsilon_{ij}\end{array} \right\}(\xi_0) =  
        \left\{ \begin{array}{lr} \hat{u}_{ij}^{\Gamma^v} \\ \tilde{u}_{ijk}^{\Gamma^v} \end{array} \right\}(\xi_0) p_j -
        \left\{ \begin{array}{lr} \hat{p}_{ij}^{\Gamma^v} \\ \tilde{p}_{ijk}^{\Gamma^v} \end{array} \right\}(\xi_0) u_j + \\
        + \left\{ \begin{array}{lr} \hat{u}_{ijk}^{\Gamma_3^w} \\ \tilde{u}_{ijkl}^{\Gamma_3^w} \end{array} \right\}(\xi_0) 
        \varepsilon_{jk}  
        - \left\{ \begin{array}{lr} \hat{u}_{ijk}^{\Omega_3^w} \\ \tilde{u}_{ijkl}^{\Omega_3^w} \end{array} \right\}(\xi_0) 
        \varepsilon_{jk},  
    \end{split}
\end{equation}

Отримаємо наступну систему рівнянь:

\begin{equation*}
    \begin{bmatrix}
        H_R^L & H_c^L & 0 & 0 & 0 \\ 
        0 & H_c^P & H_R^P & H_D^P & H_I^P
    \end{bmatrix}
    \begin{bmatrix}
        U_R^L \\ 
        U_c \\
        U_R^P \\ 
        U_D^P \\ 
        U_I^P  
    \end{bmatrix} 
    = 
    \begin{bmatrix}
        G_R^LD-G_R^LE & G_c^L & 0 & 0 & 0\\ 
        0 & -G_c^P & G_D^P & 0 & 0
    \end{bmatrix}
    \begin{bmatrix}
        U_R^L \\
        Q_c \\
        Q_D^P \\
        1 \\
        1  
    \end{bmatrix} 
\end{equation*}



\subsection{Розв'язок}

Для різних діаметрів злоякісних включень маємо наступні результати

\begin{figure}[ht!]
    \centering
    \begin{tikzpicture}
        \begin{axis}[xmin=0.19, xmax=0.21]
        \addplot table [x=x, y=y_d_0025_c_0, col sep=comma, mark=none] {data/elasticity_2d_border_data.csv};
        \addplot table [x=x, y=y_d_005_c_0, col sep=comma, mark=none] {data/elasticity_2d_border_data.csv};
        \end{axis}
    \end{tikzpicture}
    \caption{Тиск на поверхні м'якої тканини залежно від розмірів включення}
    \label{fig:elasticity_2d_solution_border_pressure_size}
\end{figure}

\begin{figure}[ht!]
    \centering
    \begin{tikzpicture}
        \begin{axis}[xmin=0.19, xmax=0.21]
        \addplot table [x=x, y=y_d_0025_c_0, col sep=comma, mark=none] {data/elasticity_2d_border_data.csv};
        \addplot table [x=x, y=y_d_0025_c_001, col sep=comma, mark=none] {data/elasticity_2d_border_data.csv};
        \end{axis}
    \end{tikzpicture}
    \caption{Тиск на поверхні м'якої тканини залежно від відстані до поверхні}
    \label{fig:elasticity_2d_solution_border_pressure_depth}
\end{figure}

Як бачимо на графіках (рис. \ref{fig:elasticity_2d_solution_border_pressure_size}) і (рис.
\ref{fig:elasticity_2d_solution_border_pressure_depth}) температура на межі двох середовищ збільшується якщо
збільшується розмір злоякісного включення $\Omega_3$ або зменшується його відстань до межі контакту.

\section{Просторова стаціонарна задача}

\subsection{Формулювання задачі}

\subsection{Дискретизація рівнянь}

Застосуємо підхід описаний у статті \cite{personal-bem-elasticity}, який полягає у поєднанні МГЕ і МСЕ. Для цього
рівняння перепишемо таким чином, щоб коефіцієнт теплопровідності був сталим майже всюди крім доданка визначеного в
області неоднорідності $\Omega_3$: 

Застосуємо закон Гука \ref{eqn:elasticity_hooke_law} і співвідношеннями Коші \ref{eqn:elasticity_gauchy_tensor} щоб 
отримати інтегральні представлення для тензора напружень, тензора деформацій і напружень:


\begin{equation}
    \label{eqn:elasticity_3d_cobem_direct_elasticity_components}
    \begin{split}
        \left\{ \begin{array}{lr} \varepsilon_{ij} \\ \sigma_{ij} \\ p_{ij} \end{array} \right \} = 5
    \end{split}
\end{equation}

\subsection{Розв'язок}

\section{Висновок}

Проведені дослідження показали, що МПГЕ забезпечує вищу\emph{ }точність розрахунків напружено-деформованого стану у
порівнянні з МГЕ при використанні однакової кількості елементів та однакового ступеня апроксимації невідомих функцій
``фіктивних'' масових сил. Це обґрунтовується тим, що пригранична область згладжує вплив введених у ній функцій.

Точність обчислень при знаходженні напружено-деформованого стану залежить від вдалого вибору форми приграничних
елементів та значення висоти \[h{}\] області \[G^{m}{}\], натомість зміна шаблону інтегрування суттєвого значення не
мала. Рекомендації щодо вибору параметра \[h{}\] можуть бути наступними: для модельних задач з відомими аналітичними
розв'язками варто виходити з апостеріорних оцінок отриманих числових результатів, для інших -- порівнювати результати
обчислень шуканих функцій у певних ділянках краю із заданими на них граничними умовами, а також точність задоволення
умов контакту на лініях контакту.

Відзначимо, що зростання (спадання) модуля зсуву в області
\[\Omega_{2}{}\] спричинює вищі (нижчі) значення компоненти
\[u_{2}^{(G)}{}\], а зростання (спадання) коефіцієнта Пуассона -- менші (більші) значення цієї компоненти при
дослідженні кусково-однорідного тіла.

Порівняльний аналіз застосування МГЕ та МПГЕ при знаходженні напружено-деформованого стану у кусково-однорідних тілах
свідчить про можливість поєднання цих методів, зокрема, на краю області варто використовувати приграничні елементи, а на
лінії контакту -- можна граничні. При цьому слід пам'ятати, що тоді необхідно створювати одне програмне забезпечення
(ПЗ) для інтегрування по внутрішніх і приграничних елементах, а інше -- по граничних, а також те, що при обчисленні
поверхневих зусиль МГЕ з'являються інтеграли в сенсі Коші, а це вимагає попереднього аналітичного виділення особливостей
(головного значення). Тому дослідникам, які прагнуть уніфікувати своє ПЗ та бажають обмежитись тільки числовим
інтегрування краще користуватись тільки МПГЕ.
